\section{Pooled conditional thermometry}
\label{sec:thermometry}

\subsection{RBM wavefunction and sampling model}
\label{sec:setup}

We use Pauli operators $\sigma_i^\mu$, $\mu\in\{x,y,z\}$, and computational-basis configurations $v=(v_1,\dots,v_N)$, $v_i=\pm1$. The transverse-field Ising chain,
\begin{equation}
H_{\mathrm I}=-J\sum_{i=1}^N\sigma_i^z\sigma_{i+1}^z-\Gamma\sum_{i=1}^N\sigma_i^x,
\label{eq:ising}
\end{equation}
with $J=1$ and periodic boundaries, supplies an exactly solvable reference energy for even $N$~\cite{Lieb1961,He2017} against which the validation procedure of Sec.~\ref{sec:validation} is checked.

An RBM with $N$ visible spins $v$, $M$ hidden spins $h$, biases $(a,b)$, and couplings $W$ defines the auxiliary energy $\mathcal E_\theta(v,h)=\sum_i a_iv_i-\sum_jb_jh_j-\sum_{ij}W_{ij}v_ih_j$. Writing $\phi_j(v)=b_j+\sum_iW_{ij}v_i$, the marginal amplitude, Born distribution, and hidden conditional means are
\begin{align}
A_\theta(v)&=e^{-\frac12\sum_ia_iv_i}\prod_{j=1}^M[2\cosh\phi_j(v)]^{1/2},\label{eq:amplitude}\\
\pi_\theta(v)&=A_\theta(v)^2/Z_\theta,\qquad \psi_\theta(v)=s(v)\sqrt{\pi_\theta(v)},\label{eq:born}\\
\mathbb E[h_j\mid v]&=\tanh\phi_j(v).\label{eq:condmean}
\end{align}
The sign $s(v)$ is fixed (Sec.~\ref{sec:intro}); it plays no role in this section. Parameters are updated with stochastic reconfiguration (SR) using the local energy $\Eloc(v)$ and the logarithmic derivatives of $A_\theta$~\cite{Sorella1998,Sorella2001}, exactly as in prior work on this system; we do not modify the optimizer and record it here only to fix notation for Sec.~\ref{sec:resources}.

A quantum annealer used as a sampler returns joint configurations $(v^{(r)},h^{(r)})_{r=1}^S$ at each SR iteration. If these were exact samples from the Gibbs distribution $q(v,h)\propto e^{-\betahw\mathcal E_\theta(v,h)}$ at some device inverse temperature $\betahw$, rescaling the submitted coefficients by $1/\beta_x$ would give an effective inverse temperature $\betaeff=\betahw/\beta_x$, and the corresponding visible marginal would be
\begin{equation}
\pi_{\theta,\beta}(v)=\frac{e^{-\beta a\cdot v}\prod_j2\cosh[\beta\phi_j(v)]}{Z_{\theta,\beta}}.
\label{eq:tempered}
\end{equation}
Equation~\eqref{eq:tempered} is a model to be tested against the sampler's actual output, not an identity that quantum annealing satisfies by construction.

\subsection{Least-squares pooled CEM}
\label{sec:cem-ls}

The pooled least-squares estimator reuses the $(v^{(r)},h^{(r)})$ pairs already produced by SR -- no fixed-condition re-query of the device is needed -- and fits
\begin{equation}
\widehat\beta_{\mathrm{LS}}=\operatorname*{arg\,min}_{\beta>0}\;
\mathcal L(\beta)=\sum_{r=1}^S\sum_{j=1}^M\big[h_j^{(r)}-\tanh(\beta\phi_j(v^{(r)}))\big]^2,
\label{eq:cem-ls}
\end{equation}
minimized by bounded scalar search over $\beta\in[0.01,50]$ in the implementation we audit. This is the estimator used in the source report and is the one for which Sec.~\ref{sec:calibration} reports residual visible-distribution error.

\subsection{What agreement of hidden conditionals does not imply}
\label{sec:cem-limits}

Write $m_j(v)=\mathbb E_q[h_j\mid v]$ for the sampler's true hidden conditional mean. The expected per-configuration loss of Eq.~\eqref{eq:cem-ls} decomposes as
\begin{equation}
\frac{\mathbb E_q[\mathcal L(\beta)]}{S}=\sum_j\mathbb E_{q(v)}\Big[\operatorname{Var}_q(h_j\mid v)+[m_j(v)-\tanh(\beta\phi_j(v))]^2\Big].
\label{eq:conditionalproof}
\end{equation}
If the hidden conditionals are exactly correct at some $\beta_0$, that value minimizes Eq.~\eqref{eq:conditionalproof} \emph{regardless of the visible marginal $q(v)$}. A minimal counterexample makes this concrete: one visible and one hidden spin, $a=0$, $b=0.2$, $W=0.4$, target $\beta=1$, so the two conditional fields are $\mp0.2,\allowbreak 0.6$ depending on $v=\mp1$. A sampler that places all visible weight on $v=-1$ but draws $h$ from the correct conditional at that $v$ reproduces $\widehat\beta_{\mathrm{LS}}=1$ exactly, while its visible total-variation error is
\begin{equation}
\DTV(q(v),\pi_\theta(v))=\frac{\cosh(0.6)}{\cosh(-0.2)+\cosh(0.6)}\approx0.5375.
\label{eq:counterexample}
\end{equation}
Any pooled-CEM estimate is therefore a diagnostic of the hidden-layer response scale, not a certificate of the full visible distribution; Sec.~\ref{sec:calibration} reports the residual visible-distribution error directly rather than inferring it from $\widehat\beta$ alone.
